\documentclass[10pt,a4paper]{amsart}
\usepackage{amsmath,amssymb,amsthm,graphicx}
% Retained minimal preamble for the Stacks Project chapter sites-cohomology.tex (Cohomology on Sites).
% Verbatim excerpt of the chapter preamble: Xy-pic setup (native preamble lines 14-18) and the native
% theorem declarations with the mathematical macros (native preamble lines 154-217). Copied mechanically;
% nothing added, removed or redefined. The wrapper class and the amsmath/amssymb/amsthm/graphicx base
% packages are supplied by the assembler, so the class/externaldocument part of the native preamble is omitted.
% For commutative diagrams we use Xy-pic
\usepackage[all]{xy}

% We use 2cell for 2-commutative diagrams.
\xyoption{2cell}
% Theorem environments.
%
\theoremstyle{plain}
\newtheorem{theorem}[subsection]{Theorem}
\newtheorem{proposition}[subsection]{Proposition}
\newtheorem{lemma}[subsection]{Lemma}

\theoremstyle{definition}
\newtheorem{definition}[subsection]{Definition}
\newtheorem{example}[subsection]{Example}
\newtheorem{exercise}[subsection]{Exercise}
\newtheorem{situation}[subsection]{Situation}

\theoremstyle{remark}
\newtheorem{remark}[subsection]{Remark}
\newtheorem{remarks}[subsection]{Remarks}

\numberwithin{equation}{subsection}

% Macros
%
\def\lim{\mathop{\mathrm{lim}}\nolimits}
\def\colim{\mathop{\mathrm{colim}}\nolimits}
\def\Spec{\mathop{\mathrm{Spec}}}
\def\Hom{\mathop{\mathrm{Hom}}\nolimits}
\def\Ext{\mathop{\mathrm{Ext}}\nolimits}
\def\SheafHom{\mathop{\mathcal{H}\!\mathit{om}}\nolimits}
\def\SheafExt{\mathop{\mathcal{E}\!\mathit{xt}}\nolimits}
\def\Sch{\mathit{Sch}}
\def\Mor{\mathop{\mathrm{Mor}}\nolimits}
\def\Ob{\mathop{\mathrm{Ob}}\nolimits}
\def\Sh{\mathop{\mathit{Sh}}\nolimits}
\def\NL{\mathop{N\!L}\nolimits}
\def\CH{\mathop{\mathrm{CH}}\nolimits}
\def\proetale{{pro\text{-}\acute{e}tale}}
\def\etale{{\acute{e}tale}}
\def\QCoh{\mathit{QCoh}}
\def\Ker{\mathop{\mathrm{Ker}}}
\def\Im{\mathop{\mathrm{Im}}}
\def\Coker{\mathop{\mathrm{Coker}}}
\def\Coim{\mathop{\mathrm{Coim}}}

% Boxtimes
%
\DeclareMathSymbol{\boxtimes}{\mathbin}{AMSa}{"02}

%
% Macros for moduli stacks/spaces
%
\def\QCohstack{\mathcal{QC}\!\mathit{oh}}
\def\Cohstack{\mathcal{C}\!\mathit{oh}}
\def\Spacesstack{\mathcal{S}\!\mathit{paces}}
\def\Quotfunctor{\mathrm{Quot}}
\def\Hilbfunctor{\mathrm{Hilb}}
\def\Curvesstack{\mathcal{C}\!\mathit{urves}}
\def\Polarizedstack{\mathcal{P}\!\mathit{olarized}}
\def\Complexesstack{\mathcal{C}\!\mathit{omplexes}}
% \Pic is the operator that assigns to X its picard group, usage \Pic(X)
% \Picardstack_{X/B} denotes the Picard stack of X over B
% \Picardfunctor_{X/B} denotes the Picard functor of X over B
\def\Pic{\mathop{\mathrm{Pic}}\nolimits}
\def\Picardstack{\mathcal{P}\!\mathit{ic}}
\def\Picardfunctor{\mathrm{Pic}}
\def\Deformationcategory{\mathcal{D}\!\mathit{ef}}
\usepackage{cleveref}
\crefname{originallemma}{original lemma}{original lemmas}
\Crefname{originallemma}{Original lemma}{Original lemmas}
\makeatletter
\newlabel{sites-lemma-sheaf-coequalizer-representable}{{00WS}{source}{}{sourcepointer-1}{}}
\newlabel{sites-lemma-sheaf-coequalizer-representable@cref}{{[originallemma][0][]00WS}{[1][0][]source}{}{}{}}
\makeatother
\crefname{originallemma}{original lemma}{original lemmas}
\Crefname{originallemma}{Original lemma}{Original lemmas}
\makeatletter
\newlabel{lemma-cech-to-cohomology-sheaf-sets}{{079Z}{source}{}{sourcepointer-2}{}}
\newlabel{lemma-cech-to-cohomology-sheaf-sets@cref}{{[originallemma][0][]079Z}{[1][0][]source}{}{}{}}
\makeatother
\title{A criterion for total acyclicity of an abelian sheaf}
\author{The Stacks Project Authors}
\date{stacks-project:sites-cohomology.tex:local-snapshot-20261011; GNU FDL 1.2 or later}
\begin{document}
\maketitle
\noindent\textit{Editorial note.} All delivered mathematical text is the author's own, copied from the pinned Stacks Project chapter file sites-cohomology.tex; no statement or proof step was rewritten, reordered or removed.\par\medskip
\noindent\textit{Editorial note.} The proof is the author's complete argument for the implication, and the statement records that the converse holds as well; the induction hypothesis and its step are both written out.\par\medskip
\noindent\textit{Editorial note.} Every labelled reference occurring inside the counted spans is carried in assembly.external\_references as a non-load-bearing navigation pointer (same-chapter ones at their definition site, cross-chapter ones at the invocation inside the counted span), so no load-bearing mathematics is replaced by a stub. No bibliography citation occurs inside the counted spans.\par\medskip
\noindent\textit{Editorial note.} The checked chapter list for same-name collisions is the set of stacks chapters staged in this repository: chow, cohomology, divisors, etale-cohomology, flat, more-algebra, more-morphisms, morphisms, spaces-more-morphisms, varieties together with this card's chapter sites-cohomology.\par\medskip
\noindent\textit{Editorial note.} The preamble is the chapter's own retained minimal preamble (sources/stacks-sites-cohomology/retained-preamble.tex), a verbatim excerpt of the native preamble: the Xy-pic setup and the native theorem declarations with the mathematical macros.\par\medskip
\noindent\textit{Allowed standard prerequisites (editorial).}\begin{itemize}
\item Totally acyclic sheaf and the extended Cech complex. Same chapter, the definition of total acyclicity and the extended Cech complex of a surjection of sheaves of sets. Supplies the conclusion of the statement and the complex whose exactness appears in hypothesis (2).
\item Vanishing of cohomology of a sheaf over a site. Same chapter, the computation of H\textasciicircum{}p(U, F) and the interpretation of the extended Cech complex as the complex computing cohomology of F over hypercoverings. Supplies hypothesis (1) and the inductive machinery of the proof.
\end{itemize}\medskip
\ifdefined\hypertarget\hypertarget{sourcepointer-1}{}\fi
\noindent\textit{Source reference.} Original lemma 00WS: sources/stacks-sites-cohomology/pinned-original-source.tex, source lines 2169--2169. Navigation pointer invoked at proof line 2169; the counted span contains the pointer only and the cited result is not reproduced inside it, so the pointer is registered as a non-load-bearing navigation stub in the established form for Stacks material. No bibliography citation is involved.\par\medskip
\ifdefined\hypertarget\hypertarget{sourcepointer-2}{}\fi
\noindent\textit{Source reference.} Original lemma 079Z: sources/stacks-sites-cohomology/pinned-original-source.tex, source lines 2073--2073. Navigation pointer invoked at proof line 2190; the counted span contains the pointer only and the cited result is not reproduced inside it, so the pointer is registered as a non-load-bearing navigation stub in the established form for Stacks material. No bibliography citation is involved.\par\medskip
\begin{lemma}
\label{lemma-characterize-limp}
Let $\mathcal{C}$ be a site. Let $\mathcal{F}$ be an abelian sheaf. If
\begin{enumerate}
\item $H^p(U, \mathcal{F}) = 0$ for $p > 0$ and $U \in \Ob(\mathcal{C})$, and
\item for every surjection $K' \to K$ of sheaves of sets the
extended {\v C}ech complex
$$
0 \to H^0(K, \mathcal{F}) \to H^0(K', \mathcal{F}) \to
H^0(K' \times_K K', \mathcal{F}) \to \ldots
$$
is exact,
\end{enumerate}
then $\mathcal{F}$ is totally acyclic (and the converse holds too).
\end{lemma}

\begin{proof}
By assumption (1) we have $H^p(h_U^\#, g^{-1}\mathcal{I}) = 0$ for all
$p > 0$ and all objects $U$ of $\mathcal{C}$. Note that if
$K = \coprod K_i$ is a coproduct of sheaves of sets on $\mathcal{C}$
then $H^p(K, g^{-1}\mathcal{I}) = \prod H^p(K_i, g^{-1}\mathcal{I})$.
For any sheaf of sets $K$ there exists a surjection
$$
K' = \coprod h_{U_i}^\# \longrightarrow K
$$
see Sites, Lemma \ref{sites-lemma-sheaf-coequalizer-representable}.
Thus we conclude that: (*) for every sheaf of sets $K$ there exists a
surjection $K' \to K$ of sheaves of sets such that $H^p(K', \mathcal{F}) = 0$
for $p > 0$. We claim that (*) and condition (2) imply that $\mathcal{F}$
is totally acyclic.
Note that conditions (*) and (2) only depend on $\mathcal{F}$ as an
object of the topos $\Sh(\mathcal{C})$ and not on the underlying site.
(We will not use property (1) in the rest of the proof.)

\medskip\noindent
We are going to prove by induction on $n \geq 0$ that (*) and (2)
imply the following induction hypothesis $IH_n$:
$H^p(K, \mathcal{F}) = 0$ for all $0 < p \leq n$ and
all sheaves of sets $K$. Note that $IH_0$ holds. Assume $IH_n$. Pick
a sheaf of sets $K$. Pick a surjection $K' \to K$ such that
$H^p(K', \mathcal{F}) = 0$ for all $p > 0$. We have a
spectral sequence with
$$
E_1^{p, q} = H^q(K'_p, \mathcal{F})
$$
covering to $H^{p + q}(K, \mathcal{F})$, see
Lemma \ref{lemma-cech-to-cohomology-sheaf-sets}.
By $IH_n$ we see that $E_1^{p, q} = 0$ for $0 < q \leq n$ and by
assumption (2) we see that $E_2^{p, 0} = 0$ for $p > 0$. Finally, we have
$E_1^{0, q} = 0$ for $q > 0$ because $H^q(K', \mathcal{F}) = 0$ by
choice of $K'$. Hence we conclude that $H^{n + 1}(K, \mathcal{F}) = 0$
because all the terms $E_2^{p, q}$ with $p + q = n + 1$ are zero.
\end{proof}

\end{document}
